% report.tex -- build with `make` (latexmk -lualatex; see the Makefile).
% Every figure, commit id and date in this file is a macro of generated/numbers.tex or a row of a
% generated/*.tex table, written by reports/programme/numbers.py from one commit (REV in the
% Makefile). Nothing here is a number typed by hand.
% preamble.tex -- the preamble of reports/status-2026-11, with: a glyph fallback for the symbols
% the programme's files carry (STIX Two Math, then DejaVu Sans), \tid{} for target ids, and the
% PDF's optional info suppressed so that two builds give the same bytes (see the Makefile).
\documentclass[11pt,a4paper]{article}

\usepackage[margin=2.2cm]{geometry}
\usepackage{amsmath,amssymb,amsthm}
\usepackage{fontspec}
\usepackage{unicode-math}
\directlua{luaotfload.add_fallback("progfallback",
  {"STIXTwoMath-Regular.otf:mode=node;", "DejaVuSans.ttf:mode=node;"})}
\setmainfont{STIX Two Text}[RawFeature={fallback=progfallback}]
\setmathfont{STIX Two Math}
\setmonofont{DejaVu Sans Mono}[Scale=0.80, RawFeature={fallback=progfallback}]
\usepackage{booktabs,longtable,array}
\usepackage{fancyvrb}
\usepackage{luacode}
\usepackage{enumitem}
\usepackage[hidelinks]{hyperref}
\setlist{itemsep=1pt,topsep=3pt}
\setlength{\emergencystretch}{3em}

% reproducible output: no PTEX.FullBanner, no PTEX.FileName (the dates and the /ID follow
% SOURCE_DATE_EPOCH, which the Makefile sets to the commit time of REV)
\pdfvariable suppressoptionalinfo \numexpr 1 + 2 \relax

% ---------------------------------------------------------------- Lean names
% \lean{name}: typeset a Lean identifier verbatim in monospace, with break points after '_', '.',
% ':' and '/', so long names do not overflow table cells. \tid{T-...}: a target id, breakable
% after '-' as well.
\begin{luacode*}
function leanname(s, breaks)
  breaks = breaks or "_.:/"
  local buf = ""
  local prev = ""
  for _, c in utf8.codes(s) do
    local ch = utf8.char(c)
    if prev:match("^%l$") and ch:match("^%u$") then
      tex.sprint(-2, buf); buf = ""
      tex.sprint("\\penalty100\\hskip0pt\\relax{}")
    end
    buf = buf .. ch
    if breaks:find(ch, 1, true) then
      tex.sprint(-2, buf); buf = ""
      tex.sprint("\\allowbreak{}")
    end
    prev = ch
  end
  tex.sprint(-2, buf)
end
\end{luacode*}
\newcommand{\lean}[1]{\texttt{\directlua{leanname("\luaescapestring{\detokenize{#1}}")}}}
\newcommand{\tid}[1]{\texttt{\directlua{leanname("\luaescapestring{\detokenize{#1}}", "-_.:/")}}}
\newcommand{\lbl}[1]{\lean{#1}}
\newcommand{\file}[1]{\lean{#1}}

% generated by reports/programme/numbers.py -- do not edit
\newcommand{\numBinderCheckViolations}{0}
\newcommand{\numBinderErsPredictabilityLayerOneLean}{22}
\newcommand{\numBinderErsPredictabilityLayerOneTarget}{6}
\newcommand{\numBinderErsPredictabilityLayerThreeEntries}{35}
\newcommand{\numBinderErsPredictabilityLayerThreeLeanEntries}{7}
\newcommand{\numBinderErsPredictabilityLayerThreeLeanRows}{0}
\newcommand{\numBinderErsPredictabilityLayerThreeModels}{2}
\newcommand{\numBinderErsPredictabilityLayerThreeRows}{28}
\newcommand{\numBinderErsPredictabilityLayerTwoLinks}{59}
\newcommand{\numBinderErsPredictabilityLayerTwoOpenEntries}{41}
\newcommand{\numBinderErsPredictabilityLayerTwoOpenNames}{7}
\newcommand{\numBinderErsPredictabilityLayerTwoOpenNonladder}{5}
\newcommand{\numBinderErsPredictabilityLayerTwoOpenProposed}{5}
\newcommand{\numBinderErsPredictabilityLayerTwoOpenRows}{14}
\newcommand{\numBinderErsPredictabilityLayerTwoProperties}{6}
\newcommand{\numBinderErsPredictabilityLayerTwoProvided}{55}
\newcommand{\numBinderErsPredictabilityLayerTwoRows}{23}
\newcommand{\numBinderErsPredictabilityLedgerRows}{32}
\newcommand{\numBinderErsPredictabilityProposed}{6}
\newcommand{\numBinderErsPredictabilityRows}{28}
\newcommand{\numBinderSbierskiMghdLayerOneLean}{16}
\newcommand{\numBinderSbierskiMghdLayerOneTarget}{3}
\newcommand{\numBinderSbierskiMghdLayerThreeEntries}{28}
\newcommand{\numBinderSbierskiMghdLayerThreeLeanEntries}{14}
\newcommand{\numBinderSbierskiMghdLayerThreeLeanRows}{0}
\newcommand{\numBinderSbierskiMghdLayerThreeModels}{3}
\newcommand{\numBinderSbierskiMghdLayerThreeRows}{14}
\newcommand{\numBinderSbierskiMghdLayerTwoLinks}{60}
\newcommand{\numBinderSbierskiMghdLayerTwoOpenEntries}{4}
\newcommand{\numBinderSbierskiMghdLayerTwoOpenNames}{1}
\newcommand{\numBinderSbierskiMghdLayerTwoOpenNonladder}{1}
\newcommand{\numBinderSbierskiMghdLayerTwoOpenProposed}{1}
\newcommand{\numBinderSbierskiMghdLayerTwoOpenRows}{4}
\newcommand{\numBinderSbierskiMghdLayerTwoProperties}{10}
\newcommand{\numBinderSbierskiMghdLayerTwoProvided}{60}
\newcommand{\numBinderSbierskiMghdLayerTwoRows}{13}
\newcommand{\numBinderSbierskiMghdLedgerRows}{19}
\newcommand{\numBinderSbierskiMghdProposed}{1}
\newcommand{\numBinderSbierskiMghdRows}{19}
\newcommand{\numBinderSeriesCommits}{10}
\newcommand{\numBinders}{2}
\newcommand{\numJournalAuthorLines}{4}
\newcommand{\numJournalEntries}{20}
\newcommand{\numJournalPlans}{5}
\newcommand{\numJournalWaves}{16}
\newcommand{\numLadderClassFails}{249}
\newcommand{\numLadderClassFailsOff}{0}
\newcommand{\numLadderClassImplied}{16}
\newcommand{\numLadderClassImpliedStanding}{0}
\newcommand{\numLadderClassOpen}{47}
\newcommand{\numLadderClassUnknown}{2}
\newcommand{\numLadderClassified}{312}
\newcommand{\numLadderConsequenceChecked}{42}
\newcommand{\numLadderConsequenceLevels}{14}
\newcommand{\numLadderConsequenceSeed}{0}
\newcommand{\numLadderConsequences}{42}
\newcommand{\numLadderFails}{68}
\newcommand{\numLadderFailsChecked}{68}
\newcommand{\numLadderFailsOff}{0}
\newcommand{\numLadderFailsSeed}{0}
\newcommand{\numLadderFirstClassified}{11}
\newcommand{\numLadderFirstConsequences}{1}
\newcommand{\numLadderFirstFails}{0}
\newcommand{\numLadderFirstImplies}{7}
\newcommand{\numLadderFirstPairs}{110}
\newcommand{\numLadderFirstProperties}{11}
\newcommand{\numLadderIllTyped}{192}
\newcommand{\numLadderImpliedStanding}{0}
\newcommand{\numLadderImplies}{25}
\newcommand{\numLadderImpliesChecked}{25}
\newcommand{\numLadderImpliesSeed}{0}
\newcommand{\numLadderOpenEntries}{47}
\newcommand{\numLadderProperties}{23}
\newcommand{\numLadderRules}{24}
\newcommand{\numLadderRulesSkipped}{1}
\newcommand{\numLadderSeriesCommits}{13}
\newcommand{\numLadderSortCausality}{18}
\newcommand{\numLadderSortGeometry}{3}
\newcommand{\numLadderSortPair}{2}
\newcommand{\numLadderTransferChecked}{7}
\newcommand{\numLadderTransferSeed}{0}
\newcommand{\numLadderTransfers}{7}
\newcommand{\numLadderViolations}{0}
\newcommand{\numLadderWithin}{314}
\newcommand{\numLedgerErsPredictabilityBucketA}{0}
\newcommand{\numLedgerErsPredictabilityBucketB}{15}
\newcommand{\numLedgerErsPredictabilityBucketC}{0}
\newcommand{\numLedgerErsPredictabilityBucketD}{7}
\newcommand{\numLedgerErsPredictabilityBucketNone}{10}
\newcommand{\numLedgerErsPredictabilityBucketT}{0}
\newcommand{\numLedgerErsPredictabilityClosed}{0}
\newcommand{\numLedgerErsPredictabilityClosedlean}{0}
\newcommand{\numLedgerErsPredictabilityFindingrows}{0}
\newcommand{\numLedgerErsPredictabilityFindings}{0}
\newcommand{\numLedgerErsPredictabilityNlproved}{23}
\newcommand{\numLedgerErsPredictabilityPartial}{22}
\newcommand{\numLedgerErsPredictabilityStated}{22}
\newcommand{\numLedgerErsPredictabilityStatements}{32}
\newcommand{\numLedgerErsPredictabilityTarget}{32}
\newcommand{\numLedgerErsPredictabilityTerminal}{0}
\newcommand{\numLedgerEttBucketA}{60}
\newcommand{\numLedgerEttBucketB}{12}
\newcommand{\numLedgerEttBucketC}{2}
\newcommand{\numLedgerEttBucketD}{59}
\newcommand{\numLedgerEttBucketNone}{3}
\newcommand{\numLedgerEttBucketT}{1}
\newcommand{\numLedgerEttClosed}{60}
\newcommand{\numLedgerEttClosedlean}{60}
\newcommand{\numLedgerEttFindingrows}{75}
\newcommand{\numLedgerEttFindings}{195}
\newcommand{\numLedgerEttNlproved}{96}
\newcommand{\numLedgerEttPartial}{71}
\newcommand{\numLedgerEttStated}{131}
\newcommand{\numLedgerEttStatements}{137}
\newcommand{\numLedgerEttTarget}{136}
\newcommand{\numLedgerEttTerminal}{1}
\newcommand{\numLedgerSbierskiMghdBucketA}{0}
\newcommand{\numLedgerSbierskiMghdBucketB}{3}
\newcommand{\numLedgerSbierskiMghdBucketC}{0}
\newcommand{\numLedgerSbierskiMghdBucketD}{16}
\newcommand{\numLedgerSbierskiMghdBucketNone}{0}
\newcommand{\numLedgerSbierskiMghdBucketT}{0}
\newcommand{\numLedgerSbierskiMghdClosed}{0}
\newcommand{\numLedgerSbierskiMghdClosedlean}{0}
\newcommand{\numLedgerSbierskiMghdFindingrows}{0}
\newcommand{\numLedgerSbierskiMghdFindings}{0}
\newcommand{\numLedgerSbierskiMghdNlproved}{17}
\newcommand{\numLedgerSbierskiMghdPartial}{19}
\newcommand{\numLedgerSbierskiMghdStated}{19}
\newcommand{\numLedgerSbierskiMghdStatements}{19}
\newcommand{\numLedgerSbierskiMghdTarget}{19}
\newcommand{\numLedgerSbierskiMghdTerminal}{0}
\newcommand{\numOqCitedDecs}{2}
\newcommand{\numOqOpen}{15}
\newcommand{\numOqRaisedDates}{2}
\newcommand{\numTargetKinds}{8}
\newcommand{\numTargets}{90}
\newcommand{\numTargetsActive}{0}
\newcommand{\numTargetsBlocked}{8}
\newcommand{\numTargetsBlockedOnOq}{8}
\newcommand{\numTargetsDone}{50}
\newcommand{\numTargetsDoneBeforeRevDate}{17}
\newcommand{\numTargetsDoneOnRevDate}{33}
\newcommand{\numTargetsDropped}{1}
\newcommand{\numTargetsNeedsRuling}{3}
\newcommand{\numTargetsPlanned}{20}
\newcommand{\numTargetsReady}{8}
\newcommand{\progJournalAuthorOqs}{OQ-69, OQ-73, OQ-77, OQ-78, OQ-80, OQ-81}
\newcommand{\progJournalFirst}{2026-11-20 19:35}
\newcommand{\progJournalLast}{2026-11-21 10:15}
\newcommand{\progLadderFirstRev}{c6dc0bb}
\newcommand{\progLadderFirstTime}{2026-11-20 18:18}
\newcommand{\progLadderToolWithinRev}{5185566}
\newcommand{\progLadderUnknownPairs}{(\lean{cauchy_slice}, \lean{rigid}), (\lean{germ_extension}, \lean{rigid})}
\newcommand{\progLadderWithinRev}{52b8b45}
\newcommand{\progLedgerErsPredictabilityLeanRev}{39be167}
\newcommand{\progLedgerEttLeanRev}{39be167}
\newcommand{\progLedgerSbierskiMghdLeanRev}{39be167}
\newcommand{\progPdfId}{59097B6E24122F2E1305EA86732A5017}
\newcommand{\progRev}{963b755}
\newcommand{\progRevDate}{21 November 2026}
\newcommand{\progRevTime}{2026-11-21 10:15}
\newcommand{\progStartRev}{c6dc0bb}
\newcommand{\progStartTime}{2026-11-20 18:18}

% the PDF's /ID: a function of REV, not of the build directory (the dates follow SOURCE_DATE_EPOCH)
\edef\progTrailerId{[<\progPdfId> <\progPdfId>]}
\pdfvariable trailerid \expandafter{\progTrailerId}
\newcommand{\gen}[1]{\csname @@input\endcsname generated/#1 }

\title{The ETT programme: progress report\\[2pt]
\large The ladder, the binders, the targets and the open questions}
\author{Generated from the programme's files at commit \texttt{\progRev}\\
\small \texttt{LeanArticles}, the \texttt{/ett-programme} loop}
\date{\progRevDate}

\begin{document}
\maketitle

\begin{abstract}
At commit \texttt{\progRev} (\progRevTime{} UTC), the causality ladder has \numLadderProperties{}
properties. Of the \numLadderWithin{} ordered pairs within one sort, \numLadderClassified{} are
classified, \numLadderClassFails{} of them by a counterexample. The \numBinders{} binders bind
\numBinderErsPredictabilityRows{} rows of \texttt{ers-predictability} and
\numBinderSbierskiMghdRows{} rows of \texttt{sbierski-mghd}. Layer~1 is in Lean for
\numBinderErsPredictabilityLayerOneLean{} and \numBinderSbierskiMghdLayerOneLean{} of these rows,
and layer~3 for \numBinderErsPredictabilityLayerThreeLeanRows{} and
\numBinderSbierskiMghdLayerThreeLeanRows{}. Of the programme's \numTargets{} targets,
\numTargetsDone{} are done (\numTargetsDoneOnRevDate{} of them on \progRevDate), \numTargetsReady{}
are ready and \numTargetsNeedsRuling{} need a ruling. \numOqOpen{} open questions of the register
mention the programme.
\end{abstract}

\tableofcontents

%======================================================================
\section{The programme}
%======================================================================

The ETT programme reads ETT as a reasoning system for causality (\file{programme/README.md}).
Every result it handles is cut into layers:
\begin{enumerate}
\item a global, causal-topological theorem of ETT, proved in any topological spacetime from
  causality properties alone and kept in \file{lean-causal/Causal/}, which imports no metric;
\item the localisation properties that theorem needs: local, checkable properties, each a node or
  an edge hypothesis of the ladder;
\item for each class of models, a geometric lemma that proves the localisation property;
\item classification criteria, which sort the cases by the causal-topological property that holds.
\end{enumerate}
Its files are the ladder (\file{programme/ladder.toml}, checked by \file{tools/ladder.py}), the
roadmap of targets (\file{programme/roadmap.toml}), a binder for each ledgered Sbierski paper
(\file{papers/<slug>/certificate/binder.toml}, checked by \file{tools/binder.py}), and the journal
(\file{programme/JOURNAL.md}), with an entry for each iteration of the loop.

Every figure below was read from commit \texttt{\progRev} by \file{reports/programme/numbers.py},
using that commit's own tools. The programme's first commit is \texttt{\progStartRev}
(\progStartTime{} UTC). This is the programme's first progress report, so each series starts at
that commit.

%======================================================================
\section{The ladder}\label{sec:ladder}
%======================================================================

\paragraph{Contents.} The ladder has \numLadderProperties{} properties, by sort:
\begin{itemize}
\item \numLadderSortCausality{} of sort \texttt{causality};
\item \numLadderSortGeometry{} of sort \texttt{geometry};
\item \numLadderSortPair{} of sort \texttt{pair}.
\end{itemize}
It records:
\begin{itemize}
\item \numLadderImplies{} implications, of which \numLadderImpliesChecked{} are checked and
  \numLadderImpliesSeed{} seed;
\item \numLadderFails{} counterexamples (\numLadderFailsChecked{} checked, \numLadderFailsSeed{}
  seed);
\item \numLadderConsequences{} consequences (\numLadderConsequenceChecked{} checked,
  \numLadderConsequenceSeed{} seed), at \numLadderConsequenceLevels{} distinct premises;
\item \numLadderTransfers{} transfers along a nesting (\numLadderTransferChecked{} checked);
\item \numLadderOpenEntries{} \texttt{[[open]]} entries.
\end{itemize}
\texttt{tools/ladder.py check} reports \numLadderViolations{} violations.

\paragraph{Classification.} \texttt{tools/ladder.py matrix} gives every ordered pair of distinct
properties a value (Table~\ref{tab:matrix}). Of the \numLadderWithin{} pairs within one sort,
\numLadderClassified{} are classified: implied, refuted by a counterexample on ETT's spacetimes, or
recorded open. A further \numLadderIllTyped{} pairs join properties of different sorts. These are \texttt{ill-typed}:
neither implication is a statement, so they are not counted. Of the rest:
\begin{itemize}
\item \numLadderClassFailsOff{} pairs are refuted only off the ambient;
\item \numLadderClassImpliedStanding{} are implied only under standing hypotheses;
\item \numLadderClassUnknown{} are unknown: \progLadderUnknownPairs.
\end{itemize}

\begin{table}[h]
\centering\small
\begin{tabular}{@{}lrrrrrrrr@{}}
\toprule
sort & properties & pairs & implied & fails & open & fails-off & implied-standing & unknown\\
\midrule
\gen{ladder-matrix}
\bottomrule
\end{tabular}
\caption{The matrix at \texttt{\progRev}, by sort. ``Pairs'' counts the ordered pairs of distinct
properties of that sort. Each pair takes the first value of the tool's order that applies.}
\label{tab:matrix}
\end{table}

\paragraph{The tactic.} From the checked edges, \texttt{tools/ladder.py emit-lean} makes
\numLadderRules{} rules of the \lean{causal_ladder} tactic. Checked edges it skips:
\numLadderRulesSkipped{}. \lean{Causal/Ladder/Rules.lean} names each skipped edge in a comment.

\paragraph{How the ladder moved.} Table~\ref{tab:ladder-series} shows the ladder at each of the
\numLadderSeriesCommits{} commits up to \texttt{\progRev} that changed \file{programme/ladder.toml}.
Each row is the output of that commit's own \texttt{tools/ladder.py check} and \texttt{matrix},
run on that commit's files. The series is therefore read from \texttt{git}, not parsed from the
journal's ladder lines, whose format varies from entry to entry.

The tool changed while the series ran: commit \texttt{\progLadderToolWithinRev} changed
\file{tools/ladder.py} so that it counts the pairs within one sort. Pair counts marked
$^\dagger$ are of all ordered pairs, as the tool counted them before that commit. The first row
read by the new tool is \texttt{\progLadderWithinRev}.

\begin{table}[h]
\centering\footnotesize
\setlength{\tabcolsep}{3pt}
\begin{tabular}{@{}ll>{\raggedright\arraybackslash}p{4.1cm}rrrrrrrrr@{}}
\toprule
commit & UTC & target & prop. & impl. & c.ex. & cons. & classified & pairs & implied & fails & open\\
\midrule
\gen{ladder-series}
\bottomrule
\end{tabular}
\caption{The ladder at each commit that changed it. ``Target'' lists the targets whose
\texttt{commit} in the roadmap is this commit. The columns implied, fails and open count the values printed by
\texttt{matrix}.}
\label{tab:ladder-series}
\end{table}

%======================================================================
\section{The binders and the Sbierski ledgers}\label{sec:binders}
%======================================================================

\paragraph{Coverage.} There are \numBinders{} binders, and \texttt{tools/binder.py check} reports
\numBinderCheckViolations{} violations. Table~\ref{tab:binders} gives, for each binder, the
counts of \texttt{tools/binder.py status} and how each layer is covered.
\begin{itemize}
\item \textbf{Layer~1} is in Lean for a row when every entry of its \texttt{ett} is a Lean name.
  Otherwise the row names a programme target.
\item \textbf{Layer~2} counts the pairs (row, ladder property) of \texttt{localisation}, and those
  that a geometric entry of the row proves. It also counts the entries of
  \texttt{open\_localisation}.
\item \textbf{Layer~3} is in Lean for a row when every geometric entry of the row cites only Lean
  names. The table also gives the geometric entries that do so.
\end{itemize}

\begin{table}[h]
\centering\small
\begin{tabular}{@{}lrrrrrrrr@{}}
\toprule
& & & \multicolumn{2}{c}{layer 1} & \multicolumn{2}{c}{layer 2} & \multicolumn{2}{c}{layer 3}\\
\cmidrule(lr){4-5}\cmidrule(lr){6-7}\cmidrule(l){8-9}
binder & ledger rows & bound & Lean & target & proved & open & rows & entries\\
\midrule
\gen{binder-coverage}
\bottomrule
\end{tabular}
\caption{Binder coverage at \texttt{\progRev}. Layer~2 ``proved'' is the pairs that a geometric
entry proves, out of all pairs. ``Open'' is the number of \texttt{open\_localisation} entries.
Layer~3 ``rows'' is the rows in Lean; ``entries'' is the geometric entries citing only Lean
names, out of all geometric entries.}
\label{tab:binders}
\end{table}

\paragraph{Open localisation properties.} A row may declare localisation properties open
(\texttt{open\_localisation}). Some of the names it uses are not ladder properties. A binder's
\texttt{[[proposed]]} entries propose such names as ladder nodes, each with the target that
would add it. Table~\ref{tab:open} counts these.

\begin{table}[h]
\centering\small
\begin{tabular}{@{}lrrrrrr@{}}
\toprule
binder & rows & entries & names & not on the ladder & of these, proposed & \texttt{[[proposed]]}\\
\midrule
\gen{binder-open}
\bottomrule
\end{tabular}
\caption{The \texttt{open\_localisation} entries at \texttt{\progRev}: the rows that have any,
the entries, the distinct names, those that are not ladder properties, and those among them that
have a \texttt{[[proposed]]} entry. The last column counts the binder's
\texttt{[[proposed]]} entries.}
\label{tab:open}
\end{table}

\begin{table}[h]
\centering\small
\begin{tabular}{@{}llrrr@{}}
\toprule
binder & model & geometric entries & Lean names only & citing a target\\
\midrule
\gen{binder-models}
\bottomrule
\end{tabular}
\caption{Layer~3 by model class at \texttt{\progRev}.}
\label{tab:models}
\end{table}

\paragraph{The ledgers.} Table~\ref{tab:ledgers} gives the ledgers of the binders' papers. It
includes the ETT ledger, whose labels the ladder's properties carry, and uses the figures that
\texttt{tools/status\_report.py} reads for \file{reports/status-2026-11}. The ledgers read
\file{lean-causal} at commits \texttt{\progLedgerErsPredictabilityLeanRev},
\texttt{\progLedgerSbierskiMghdLeanRev} and \texttt{\progLedgerEttLeanRev}.
\begin{itemize}
\item In \texttt{ers-predictability}, \numLedgerErsPredictabilityStated{} of
  \numLedgerErsPredictabilityStatements{} rows are Lean-stated:
  \numLedgerErsPredictabilityClosedlean{} closed and \numLedgerErsPredictabilityPartial{} partial.
\item In \texttt{sbierski-mghd}, \numLedgerSbierskiMghdStated{} of
  \numLedgerSbierskiMghdStatements{} rows are Lean-stated: \numLedgerSbierskiMghdClosedlean{}
  closed and \numLedgerSbierskiMghdPartial{} partial.
\end{itemize}

\begin{table}[h]
\centering\small
\begin{tabular}{@{}lrrrrrrrrrrr@{}}
\toprule
& & \multicolumn{3}{c}{Lean} & \multicolumn{6}{c}{bucket} & \\
\cmidrule(lr){3-5}\cmidrule(lr){6-11}
ledger & rows & stated & closed & partial & A & B & C & D & T & none & closure\\
\midrule
\gen{ledgers}
\bottomrule
\end{tabular}
\caption{The ledgers at \texttt{\progRev} (\texttt{tools/ledger.py}'s header and
\texttt{ledger.json}). ``None'' counts the rows that have no bucket. Closure is the number of
rows closed in bucket~A, over the non-terminal rows.}
\label{tab:ledgers}
\end{table}

\paragraph{How the binders moved.} Table~\ref{tab:binder-series} shows \texttt{tools/binder.py
status} at each of the \numBinderSeriesCommits{} commits up to \texttt{\progRev} that changed a
binder. Each row is the output of that commit's own tool, run on that commit's files.

\begin{table}[h]
\centering\footnotesize
\begin{tabular}{@{}ll>{\raggedright\arraybackslash}p{4.2cm}rrrrrr@{}}
\toprule
& & & \multicolumn{3}{c}{\texttt{ers-predictability}} & \multicolumn{3}{c}{\texttt{sbierski-mghd}}\\
\cmidrule(lr){4-6}\cmidrule(l){7-9}
commit & UTC & target & bound & layer 1 & layer 3 & bound & layer 1 & layer 3\\
\midrule
\gen{binder-series}
\bottomrule
\end{tabular}
\caption{The binders at each commit that changed one. Each count is the tool's own: rows bound,
rows with layer~1 in Lean, and rows with layer~3 in Lean. A dash means the binder did not exist
yet.}
\label{tab:binder-series}
\end{table}

%======================================================================
\section{Targets}\label{sec:targets}
%======================================================================

\paragraph{By status and kind.} \file{programme/roadmap.toml} has \numTargets{} targets:
\begin{itemize}
\item \numTargetsDone{} done;
\item \numTargetsReady{} ready;
\item \numTargetsNeedsRuling{} needing a ruling;
\item \numTargetsBlocked{} blocked on an open question;
\item \numTargetsPlanned{} planned;
\item \numTargetsActive{} active;
\item \numTargetsDropped{} dropped.
\end{itemize}
Table~\ref{tab:kinds} breaks these down by kind, the roadmap's \texttt{kind} field.

\begin{table}[h]
\centering\small
\begin{tabular}{@{}lrrrrrrrr@{}}
\toprule
kind & planned & ready & active & needs-ruling & blocked & done & dropped & all\\
\midrule
\gen{targets-kind}
\bottomrule
\end{tabular}
\caption{Targets by kind and status at \texttt{\progRev}.}
\label{tab:kinds}
\end{table}

\paragraph{Done on \progRevDate.} A target's done date is the commit time, in UTC, of the
commit that its \texttt{commit} field records. \numTargetsDoneOnRevDate{} targets were done on
\progRevDate{} (Table~\ref{tab:done}), and \numTargetsDoneBeforeRevDate{} before.

\begin{small}
\begin{longtable}{@{}>{\raggedright\arraybackslash}p{6.6cm}llr@{}}
\caption{Targets done on \progRevDate, in order of commit time.}\label{tab:done}\\
\toprule
target & kind & commit & UTC\\
\midrule
\endfirsthead
\toprule
target & kind & commit & UTC\\
\midrule
\endhead
\bottomrule
\endfoot
\gen{targets-done-today}
\end{longtable}
\end{small}

\paragraph{Ready.} Table~\ref{tab:ready} lists the ready targets in order of priority. Each title
is the first clause of the target's own title.

\begin{small}
\begin{longtable}{@{}>{\raggedright\arraybackslash}p{4.4cm}lr>{\raggedright\arraybackslash}p{7.6cm}@{}}
\caption{Ready targets at \texttt{\progRev}.}\label{tab:ready}\\
\toprule
target & kind & priority & title (first clause)\\
\midrule
\endfirsthead
\toprule
target & kind & priority & title (first clause)\\
\midrule
\endhead
\bottomrule
\endfoot
\gen{targets-ready}
\end{longtable}
\end{small}

%======================================================================
\section{Open questions for the author}\label{sec:open}
%======================================================================

\paragraph{The register's questions.} \file{registers/DECISIONS.md} has \numOqOpen{} open
questions whose row mentions the ETT programme (Table~\ref{tab:oq}). The title of each is the
first clause of its question, without the opening ``ETT programme,''. The row itself is the
authority, and this report reads nothing further into it. The table also gives:
\begin{itemize}
\item the rulings each row cites, \numOqCitedDecs{} in all (Table~\ref{tab:dec});
\item the targets that the roadmap blocks on each question, \numTargetsBlockedOnOq{} in all.
\end{itemize}

\begin{small}
\begin{longtable}{@{}lr>{\raggedright\arraybackslash}p{7.7cm}l>{\raggedright\arraybackslash}p{2.8cm}@{}}
\caption{Open questions raised by the programme at \texttt{\progRev}.}\label{tab:oq}\\
\toprule
OQ & raised & title (first clause of the row) & cites & blocked targets\\
\midrule
\endfirsthead
\toprule
OQ & raised & title (first clause of the row) & cites & blocked targets\\
\midrule
\endhead
\bottomrule
\endfoot
\gen{oq}
\end{longtable}
\end{small}

\begin{table}[h]
\centering\small
\begin{tabular}{@{}lr>{\raggedright\arraybackslash}p{11.4cm}@{}}
\toprule
DEC & date & decision (first clause of the row)\\
\midrule
\gen{dec}
\bottomrule
\end{tabular}
\caption{The rulings that the programme's open questions cite.}
\label{tab:dec}
\end{table}

\paragraph{Targets waiting on a ruling.} \numTargetsNeedsRuling{} targets need a ruling
(Table~\ref{tab:ruling}). Each target's question is in its roadmap block. The table gives the
question's first clause.

\begin{small}
\begin{longtable}{@{}>{\raggedright\arraybackslash}p{4.4cm}l>{\raggedright\arraybackslash}p{10.0cm}@{}}
\caption{Targets needing a ruling at \texttt{\progRev}.}\label{tab:ruling}\\
\toprule
target & kind & question (first clause)\\
\midrule
\endfirsthead
\toprule
target & kind & question (first clause)\\
\midrule
\endhead
\bottomrule
\endfoot
\gen{targets-ruling}
\end{longtable}
\end{small}

\paragraph{The journal.} \file{programme/JOURNAL.md} has \numJournalEntries{} entries, from
\progJournalFirst{} to \progJournalLast{} UTC. Of these, \numJournalWaves{} report a wave and
\numJournalPlans{} a planning run; an entry may report both. Its \numJournalAuthorLines{} ``for
the author'' lines name \progJournalAuthorOqs. Table~\ref{tab:journal} gives the first clause of
each line that names no OQ row.

\begin{table}[h]
\centering\small
\begin{tabular}{@{}l>{\raggedright\arraybackslash}p{12.6cm}@{}}
\toprule
entry (UTC) & line (first clause)\\
\midrule
\gen{journal-author}
\bottomrule
\end{tabular}
\caption{Journal lines for the author that name no OQ row.}
\label{tab:journal}
\end{table}

\paragraph{Provenance.} Every figure is a macro of \file{generated/numbers.tex}, and every table
body is a file of \file{generated/}. Both are written by \file{reports/programme/numbers.py}
from commit \texttt{\progRev}, extracted with \texttt{git archive}, through that commit's
\file{tools/ladder.py}, \file{tools/binder.py}, \file{tools/roadmap.py} and
\file{tools/status_report.py}. The series run each earlier commit's own tools. No model was
called and no Lean was built. \texttt{make check} in this directory does the following:
\begin{itemize}
\item it regenerates the numbers and compares them byte for byte;
\item it builds the PDF a second time and compares it byte for byte;
\item it counts the undefined references and missing glyphs in the log;
\item it checks this text and the generated tables against every legacy draft
  (\file{papers/*/legacy/orig}), with the shingles of \file{tools/no_legacy_prose.py}.
\end{itemize}

\end{document}
